\subsection*{2.3. Finiteness theorem (in equal characteristic zero)}

\subsubsection*{Theorem 2.3.1.}
\underline{Let $S$ be a henselian trait as in 0.2.5, let $\Lambda$ be
a noetherian torsion ring, and let $f:X\longrightarrow S$ be a
morphism of finite type.  Suppose that $S$ has equal characteristic
zero.  Then, for every constructible sheaf of $\Lambda$-modules $F$
on $X_\eta$, the sheaves $R^i\Psi_{\bar\eta}(F)$ on $X_{\bar s}$ are
constructible.}

We begin by repeating the opening of the proof of SGA~4, XIV, 1.10.
Apply SGA~4, IX, 2.14(ii).  There is a resolution $F^\bullet$ of $F$
\begin{center}
\begin{tikzcd}[column sep=large, row sep=large]
{} & F \arrow[d] & {} & {} \\
0 \arrow[r] & F^0 \arrow[r] & F^1 \arrow[r] & \cdots
\end{tikzcd}
\end{center}
by sheaves of the form $\bigoplus_i p'_{i*}G_i$, where
$p'_i:Y_{i\eta}\longrightarrow X_\eta$ is finite and $G_i$ is a
constructible constant sheaf.  We may suppose that $Y_{i\eta}$ is
reduced; let $p_i:Y_i\longrightarrow X$ be the normalization of $X$
in $Y_{i\eta}$.  Using the spectral sequence of the resolution
$F^\bullet$ and applying (2.1.7.1) in the elementary case of the
finite morphism $p_i$, we see that it suffices to prove the theorem
for each pair $(Y_i,G_i)$.

We next prove the theorem when $F$ is the constant sheaf defined by a
$\Lambda$-module $M$.  Recall that $\Lambda$ is a
$\mathbb Z/n$-algebra for a suitable $n$.  The module $M$ may
therefore be d\'{e}viss\'{e} into modules $M_i$ over
$\Lambda/\ell_i\Lambda$, where each $\ell_i\mid n$ is a suitable
prime.  This reduces us to the case in which $\Lambda$ is a
$\mathbb Z/\ell$-algebra, with $\ell$ prime.  Formula (2.1.13.1) then
gives
\[
M\otimes_{\mathbb Z/\ell}
R^i\Psi_{\bar\eta}(\mathbb Z/\ell)
\xrightarrow{\ \sim\ }
R^i\Psi_{\bar\eta}(M).
\]
It remains to treat the constant sheaf $F=\mathbb Z/\ell$.

Choose a descending sequence of closed subsets
\[
X_\eta=X'_0\supset\cdots\supset X'_n
\]
such that $(X'_i-X'_{i+1})_{\mathrm{red}}$ is smooth over $\eta$, and
let $X_i=\overline{X'_i}$ be the closure in $X$.  The constant sheaf
$\mathbb Z/\ell$ is d\'{e}viss\'{e} into the constant sheaves
$\mathbb Z/\ell$ on the strata $X'_i-X'_{i+1}$, extended by zero.
Thus:

\medskip
\noindent\underline{Reduction 1.}
It is enough to prove Theorem 2.3.1 for a sheaf
$j_!\mathbb Z/\ell$, where
$j:U\hookrightarrow X_\eta$ is the inclusion of a smooth open subset.

By resolution of singularities there is a morphism
$g:X_1\longrightarrow X$ such that

\medskip
\noindent(a) $g$ is proper and induces an isomorphism
$g^{-1}(U)\xrightarrow{\sim}U$;

\noindent(b) $X_1$ is regular and $X_1-g^{-1}(U)$ is a
normal-crossings divisor in $X_1$, which is a union of regular
divisors $D_\alpha$.

Let $j_1:g^{-1}(U)\hookrightarrow X_1$.  Then
\[
Rg_*j_{1!}\mathbb Z/\ell=j_!\mathbb Z/\ell.
\]
Applying (2.1.7.1) and the finiteness theorem SGA~4, XIV, 1.1, we see
that it suffices to prove Theorem 2.3.1 for
$(X_1,j_{1!}\mathbb Z/\ell)$.

There is a resolution
\[
0\longrightarrow j_{1!}\mathbb Z/\ell
\longrightarrow\mathbb Z/\ell
\longrightarrow\bigoplus_\alpha(\mathbb Z/\ell)_{D_\alpha}
\longrightarrow
\bigoplus_{\alpha,\beta}
(\mathbb Z/\ell)_{D_\alpha\cap D_\beta}
\longrightarrow\cdots,
\]
and hence:

\medskip
\noindent\underline{Reduction 2.}
It is enough to prove Theorem 2.3.1 for the constant sheaf
$\mathbb Z/\ell$ on a regular $S$-scheme $X$ such that $X_s$ is a
normal-crossings divisor in $X$.

In this final case one uses the explicit computation I, 3.3, which in
turn uses the purity theorem SGA~4, XIX, 3.2 and the characteristic-zero
hypothesis.
